% ============================================================
%  TFM — Estudio Comparativo del Espacio Latente de Modelos
%         Generativos de Texto-Imagen
%  Autor: Víctor Menéndez Ríos
% ============================================================
\documentclass[12pt, a4paper, twoside]{book}

% --- Encoding & Language -----------------------------------
\usepackage[utf8]{inputenc}
\usepackage[T1]{fontenc}
\usepackage[spanish]{babel}

% --- Typography --------------------------------------------
\usepackage{lmodern}
\usepackage{microtype}
\usepackage{setspace}
\onehalfspacing

% --- Maths -------------------------------------------------
\usepackage{amsmath, amssymb, amsthm}
\usepackage{bm}

% --- Layout ------------------------------------------------
\usepackage[top=3cm, bottom=3cm, inner=3.5cm, outer=2.5cm]{geometry}
\usepackage{emptypage}
\pagestyle{plain}

% --- Figures & Tables --------------------------------------
\usepackage{graphicx}
\graphicspath{{./}{./data/}{./imagenes_manuscrito/}}
\usepackage{booktabs}
\usepackage{multirow}
\usepackage{array}
\usepackage{longtable}
\usepackage{caption}
\usepackage{subcaption}
\usepackage{float}
\usepackage{xcolor}

% --- Links (sin color) -------------------------------------
\usepackage[hidelinks]{hyperref}
\usepackage{url}

% --- Bibliografía ------------------------------------------
\usepackage[authoryear, round, sort&compress]{natbib}

% --- Listas ------------------------------------------------
\usepackage{enumitem}

% --- Custom commands ---------------------------------------
\newcommand{\z}{\mathbf{z}}
\newcommand{\x}{\mathbf{x}}
\newcommand{\eps}{\boldsymbol{\epsilon}}
\newcommand{\R}{\mathbb{R}}
\newcommand{\E}{\mathbb{E}}
\newcommand{\norm}[1]{\left\|#1\right\|_2}
\newcommand{\TODO}[1]{\textcolor{red}{\textbf{[PENDIENTE: #1]}}}
\newcommand{\PLACEHOLDER}[1]{\textcolor{orange}{\textit{[#1]}}}

% ===========================================================
\begin{document}
% ===========================================================

% --- Portada -----------------------------------------------
\begin{titlepage}
  \centering
  \vspace*{2cm}
  {\Large Universidad Politécnica de Madrid\par}
  \vspace{0.5cm}
  {\large Máster en Deep Learning\par}
  \vspace{3cm}
  \rule{\linewidth}{0.5pt}\\[0.4cm]
  {\LARGE\bfseries Estudio Comparativo del Espacio Latente\\[0.3cm]
   de Modelos Generativos de Texto-Imagen\par}
  \rule{\linewidth}{0.5pt}\\[1.5cm]
  {\large\itshape Trabajo de Fin de Máster\par}
  \vspace{2cm}
  {\large Autor: \textbf{Victor Mendez Riosalido}\par}
  \vspace{0.5cm}
  {\large Tutor: \textbf{Silvia Alba Uribe Mayoral}\par}
  \vfill
  {\large Curso académico 2025--2026\par}
\end{titlepage}

% --- Páginas preliminares ----------------------------------
\frontmatter

\chapter*{Resumen}
\addcontentsline{toc}{chapter}{Resumen}

Los modelos de difusión latente han transformado la síntesis de imágenes a partir de texto, pero la estructura interna de su espacio latente permanece insuficientemente comprendida. Este trabajo presenta un análisis comparativo cuantitativo del espacio latente $\z_0$ de tres modelos de la familia Stable Diffusion ---SD~1.5, SD~2.1 y SDXL--- con distintos codificadores de texto y capacidades de UNet, manteniendo constante la arquitectura del autoencoder variacional. El estudio se articula en tres ejes: análisis geométrico de las distancias latentes, análisis semántico de las direcciones de atributo y análisis del impacto del fine-tuning mediante LoRA sobre un corpus de pintura clásica (WikiArt, 500 imágenes). Los resultados muestran que SDXL presenta propiedades cualitativamente distintas en pretrain: mayor separabilidad semántica, direcciones de atributo escena-independientes con consistencia intra-tipo de 0,71 para el color, y concentración del 95\% de la energía semántica en los diez primeros componentes principales. Los tres modelos exhiben una dimensionalidad intrínseca de aproximadamente doce dimensiones efectivas, con independencia de su dimensión nominal. El fine-tuning sobre WikiArt revela un patrón contrario al esperado: SDXL, con la representación pretrain más organizada, sufre la mayor perturbación geométrica y semántica ---desplazamiento latente medio 3,5 veces superior al de los otros dos modelos, alineación PCA post-tuning de 0,54 frente a 0,82--0,89, y caídas de consistencia del 22--52\% según el atributo--- mientras que SD~1.5 y SD~2.1 muestran una deriva prácticamente nula. Este resultado evidencia que una organización semántica más estructurada no protege frente a la disrupción por adaptación de dominio, sino que la amplifica.

\bigskip
\noindent\textbf{Palabras clave:} modelos de difusión latente, espacio latente, análisis semántico, fine-tuning, LoRA, Stable Diffusion, text-to-image.

\chapter*{Abstract}
\addcontentsline{toc}{chapter}{Abstract}

Latent diffusion models have transformed text-to-image synthesis, yet the internal structure of their latent space remains poorly understood. This work presents a quantitative comparative analysis of the $\z_0$ latent space of three models from the Stable Diffusion family---SD~1.5, SD~2.1, and SDXL---differentiated by their text encoder and UNet capacity while keeping the variational autoencoder architecture constant. The study is organized around three axes: geometric analysis of latent distances, semantic analysis of attribute directions, and analysis of the impact of LoRA fine-tuning on a classical painting corpus (WikiArt, 500 images). Pre-training results show that SDXL exhibits qualitatively different properties: greater semantic separability, scene-independent attribute directions with an intra-type consistency of 0.71 for color, and concentration of 95\% of semantic energy in the top ten principal components. All three models display an intrinsic dimensionality of approximately twelve effective dimensions, regardless of their nominal dimension. WikiArt fine-tuning reveals a pattern contrary to initial predictions: SDXL, with the most organized pre-training representation, suffers the largest geometric and semantic perturbation---mean latent displacement 3.5$\times$ larger than the other two models, post-tuning PCA alignment of 0.54 versus 0.82--0.89, and consistency drops of 22--52\% across attributes---while SD~1.5 and SD~2.1 show near-negligible drift. This result demonstrates that a more structured semantic organization does not protect against disruption by domain adaptation, but amplifies it.

\bigskip
\noindent\textbf{Keywords:} latent diffusion models, latent space, semantic analysis, fine-tuning, LoRA, Stable Diffusion, text-to-image.

\tableofcontents
\listoftables
\listoffigures

% -----------------------------------------------------------
\mainmatter
% -----------------------------------------------------------

% ===========================================================
%  INTRODUCCIÓN  (incluye Estado del Arte)
% ===========================================================
\chapter{Introducción}
\label{ch:intro}

\section{Contexto y motivación}
\label{sec:motivacion}

La síntesis de imágenes a partir de descripciones textuales ha experimentado una transformación sin precedentes durante los últimos años. En agosto de 2022, la publicación de Stable Diffusion como modelo de código abierto marcó un punto de inflexión: por primera vez, un sistema capaz de producir imágenes de calidad comparable a los mejores sistemas propietarios era accesible para cualquier usuario con una tarjeta gráfica de consumo \citep{rombach2022ldm}. Desde entonces, el ecosistema ha crecido a un ritmo excepcional: modelos como SDXL \citep{podell2023sdxl}, Stable Diffusion 3 \citep{esser2024sd3} y Flux \citep{bfl2024flux} han elevado sucesivamente el techo de calidad en el ámbito del código abierto, mientras que sistemas propietarios como DALL-E 3 \citep{betker2023dalle3}, Midjourney v6 o Adobe Firefly han redefinido los estándares de usabilidad y fidelidad visual en entornos comerciales.

Sin embargo, a pesar de la impresionante calidad de las imágenes generadas, la comprensión profunda de los mecanismos internos que sustentan la síntesis permanece notablemente incompleta. En particular, el espacio latente, que se define como la representación comprimida y abstracta del contenido visual que los modelos construyen durante el entrenamiento, es la pieza central que determina la calidad, la coherencia semántica y la controlabilidad de la generación. La estructura geométrica y semántica de este espacio está lejos de ser completamente comprendida, especialmente en el contexto de los modelos de difusión latente, donde la representación latente resulta menos intuitiva que en arquitecturas anteriores como las redes generativas adversariales \citep{kwon2023diffusion}. Esta opacidad tiene consecuencias prácticas directas: las ediciones basadas en manipulación de latentes pueden producir artefactos impredecibles, las interpolaciones no siempre son semánticamente suaves, y el comportamiento de los modelos ante el fine-tuning sobre dominios específicos es difícil de anticipar.

El presente trabajo aborda esta brecha desde una perspectiva comparativa y cuantitativa. En lugar de evaluar los modelos por la calidad visual de sus salidas lo cual ya extensamente tratado en la literatura, se propone analizar y contrastar la estructura interna del espacio latente de tres modelos de difusión latente de código abierto con arquitecturas deliberadamente distintas: Stable Diffusion 1.5, Stable Diffusion 2.1 y Stable Diffusion XL. Los resultados son relevantes tanto para la comunidad investigadora que necesita referencias cuantitativas para comprender cómo la arquitectura afecta a la organización de las representaciones internas como para los profesionales que desean maximizar la controlabilidad y la predictibilidad de las ediciones.

\section{Perspectiva histórica de la síntesis visual computacional}
\label{sec:historia}

\subsection{Generación procedural y estadística (hasta 2012)}

Los primeros sistemas de síntesis visual computacional eran enteramente procedurales: algoritmos deterministas que producían imágenes a partir de reglas matemáticas explícitas. El ruido de Perlin \citep{perlin1985noise}, introducido para la generación de texturas naturales en producción cinematográfica, es un ejemplo paradigmático. Los métodos estadísticos de síntesis de texturas de \citet{efros1999texture} demostraron que era posible sintetizar nuevas texturas a partir de un ejemplo mediante muestreo no paramétrico. El denominador común de todos estos enfoques es que la representación del mundo visual estaba codificada de forma explícita y manual, lo que limitaba drásticamente la capacidad de generalización.

\subsection{El giro neural: representaciones aprendidas (2012--2014)}

La publicación de AlexNet \citep{krizhevsky2012alexnet} abrió la puerta a un nuevo paradigma: en lugar de especificar qué buscar en una imagen, los sistemas aprenderían qué representaciones son útiles directamente de los datos. La transferencia de estilo neuronal \citep{gatys2016style} demostró que las representaciones profundas contenían información visual estructurada y manipulable. El autoencoder variacional (VAE) de \citet{kingma2013vae} formalizó el primer modelo generativo profundo con un espacio latente continuo y semánticamente estructurado que podía muestrearse de forma controlada.

\subsection{La primera generación de modelos generativos profundos (2013--2017)}

Cuatro familias de modelos generativos profundos surgieron durante este período. Los autoencoders variacionales \citep{kingma2013vae} aprendían un espacio latente continuo maximizando la cota inferior de la log-verosimilitud. Las redes generativas adversariales \citep{goodfellow2014gans} adoptaron un juego de suma cero entre generador y discriminador; DCGAN \citep{radford2015dcgan} demostró que el espacio latente tenía propiedades aritméticas sorprendentes: la operación $\z_{\text{hombre con gafas}} - \z_{\text{hombre}} + \z_{\text{mujer}} \approx \z_{\text{mujer con gafas}}$ funcionaba de forma consistente. Los modelos autorregresivos \citep{vandeoord2016pixelrnn} y los flujos normalizantes \citep{dinh2017realnvp,kingma2018glow} completaron el panorama generativo de la época.

\subsection{La era GAN: síntesis de alta fidelidad (2017--2021)}

StyleGAN \citep{karras2019stylegan} y StyleGAN2 \citep{karras2020stylegan2} introdujeron la red de mapeo $f: \z \rightarrow \mathbf{w}$ que transforma el ruido latente en un espacio intermedio $\mathcal{W}$ con disentanglement superior. Trabajos como InterFaceGAN \citep{shen2020interfacegan} y GANSpace \citep{harkonen2020ganspace} demostraron que atributos semánticos complejos (e.g. edad, género, expresión, pose o iluminación) corresponden a hiperplanos lineales o componentes principales bien definidos en $\mathcal{W}$. Esta capacidad de análisis y edición lineal en el espacio latente es el principal antecedente metodológico del presente trabajo.

\subsection{El pivote hacia la difusión y la generación guiada por texto (2021--2022)}

CLIP \citep{radford2021clip} entrenó, mediante aprendizaje contrastivo a escala masiva sobre 400 millones de pares imagen-texto, un espacio de embedding compartido entre modalidades visuales y lingüísticas. \citet{rombach2022ldm} introdujeron los modelos de difusión latente (LDM), desplazando el proceso de difusión del espacio de píxeles al espacio latente comprimido de un VAE. Stable Diffusion, su implementación de código abierto entrenada sobre LAION-5B \citep{schuhmann2022laion}, se convirtió en el primer sistema de generación texto-imagen de alta calidad accesible sin infraestructura computacional extraordinaria.

\subsection{Consolidación y madurez del ecosistema (2022--2025)}

SDXL \citep{podell2023sdxl} introdujo un UNet de 2.600 millones de parámetros con doble codificador de texto y resolución nativa de $1024 \times 1024$ píxeles. Stable Diffusion 3 \citep{esser2024sd3} supuso un cambio de paradigma hacia transformadores de difusión multimodales y flow matching rectificado \citep{lipman2022flowmatching,liu2022rectified}. Flux \citep{bfl2024flux}, con 12.000 millones de parámetros, estableció el nuevo estado del arte en texto-imagen de código abierto a mediados de 2024. En el ámbito propietario, DALL-E 3 \citep{betker2023dalle3} demostró que la estrategia de \textit{recaptioning} con modelos de lenguaje mejora drásticamente la fidelidad textual sin cambios arquitectónicos.

\section{Fundamentos de los modelos generativos profundos}
\label{sec:fundamentos}

\subsection{Autoencoders variacionales}

El autoencoder variacional \citep{kingma2013vae} aprende simultáneamente un codificador probabilístico $q_\phi(\z \mid \x)$ y un decodificador $p_\theta(\x \mid \z)$. La función objetivo es la \textit{cota inferior de evidencia} (ELBO):

\begin{equation}
\mathcal{L}_{\text{ELBO}} = \E_{q_\phi(\z \mid \x)}\!\left[\log p_\theta(\x \mid \z)\right]
  - D_{\mathrm{KL}}\!\left[q_\phi(\z \mid \x) \,\|\, p(\z)\right]
\label{eq:elbo}
\end{equation}

La regularización mediante la divergencia KL impone una estructura global que hace el espacio latente interpolable: puntos cercanos en $\mathcal{Z}$ corresponden a imágenes perceptualmente similares.

\subsection{Redes generativas adversariales}

Las GANs \citep{goodfellow2014gans} definen un juego de suma cero entre generador $G: \z \rightarrow \x$ y discriminador $D: \x \rightarrow [0,1]$:

\begin{equation}
\min_G \max_D \;\E_{\x \sim p_{\text{data}}}\!\left[\log D(\x)\right]
+ \E_{\z \sim p(\z)}\!\left[\log\!\left(1 - D(G(\z))\right)\right]
\end{equation}

La relevancia de las GANs para este trabajo es metodológica: gran parte del instrumental de análisis del espacio latente, haciendo referencia a la búsqueda de direcciones semánticas lineales, PCA de activaciones, clasificadores lineales, fue desarrollado inicialmente para GANs y se traslada directamente al estudio de los LDMs.

\subsection{Modelos de difusión probabilísticos}

Los modelos de difusión probabilísticos \citep{ho2020ddpm} definen un proceso hacia adelante que añade ruido gaussiano en $T$ pasos:

\begin{equation}
q(\x_t \mid \x_0) = \mathcal{N}\!\left(\sqrt{\bar{\alpha}_t}\,\x_0,\; (1 - \bar{\alpha}_t)\,\mathbf{I}\right),
\quad \bar{\alpha}_t = \prod_{s=1}^t (1 - \beta_s)
\end{equation}

La red neuronal $\eps_\theta$ se entrena para predecir el ruido añadido:

\begin{equation}
\mathcal{L}_{\text{simple}} = \E_{\x_0,\,\eps \sim \mathcal{N}(\mathbf{0},\mathbf{I}),\,t}\!
\left[\left\|\eps - \eps_\theta\!\left(\sqrt{\bar{\alpha}_t}\,\x_0 + \sqrt{1 - \bar{\alpha}_t}\,\eps,\;t\right)\right\|_2^2\right]
\end{equation}

El muestreo DDIM \citep{song2021ddim} permite recorrer la trayectoria en $S \ll T$ pasos con naturaleza determinista, habilitando la inversión de imágenes reales.
\begin{figure}[htb]
\centering
\includegraphics[width=\textwidth]{X-Y_plot_of_algorithmically-generated_AI_art_of_European-style_castle_in_Japan_demonstrating_DDIM_diffusion_steps}
\caption{Diagrama del proceso de denoising DDIM para distintas configuraciones de pasos de inferencia y escala de guiado. La naturaleza determinista del muestreo DDIM con semilla fija garantiza la reproducibilidad del espacio latente $\z_0$ utilizado en este trabajo. Fuente: Benlisquare (2022), Wikimedia Commons. Licencia: CC~BY-SA~4.0.}
\label{fig:ddim_steps}
\end{figure}
El guiado clasificador-libre \citep{ho2022cfg} interpola entre las predicciones de ruido condicionada e incondicionada:

\begin{equation}
\hat{\eps}_\theta(\x_t, t, y) = \eps_\theta(\x_t, t, \emptyset)
+ w\bigl[\eps_\theta(\x_t, t, y) - \eps_\theta(\x_t, t, \emptyset)\bigr]
\end{equation}

donde $w = 7{,}5$ en todos los experimentos de este trabajo.

\subsection{Modelos de difusión latente}

\begin{figure}[htb]
\centering
\includegraphics[width=0.88\textwidth]{Stable_Diffusion_architecture}
\caption{Arquitectura de un modelo de difusión latente (LDM). El codificador perceptual $\mathcal{E}$ comprime la imagen $\x$ al espacio latente $\z_0$; el UNet condicionado $\eps_\theta$ aprende a revertir el proceso de difusión; y el decodificador $\mathcal{D}$ reconstruye la imagen a resolución original. Adaptado de \citet{rombach2022ldm}.}
\label{fig:sd_architecture}
\end{figure}

\citet{rombach2022ldm} propusieron desplazar el proceso de difusión al espacio latente comprimido de un VAE. Para un factor de compresión $f = 8$ y $C = 4$ canales latentes, una imagen de $512 \times 512$ se comprime en $4 \times 64 \times 64$, un factor de reducción dimensional de 192. El espacio latente denoised final antes de la decodificación o lantente $\z_0$ es el objeto central de análisis en este trabajo. Sus propiedades geométricas y semánticas reflejan directamente cómo el modelo ha organizado el contenido visual durante el proceso de difusión.

\subsection{Transformadores de difusión y flow matching}

DiT \citep{peebles2023dit} propuso sustituir el UNet por un transformer que opera sobre parches del latente, siguiendo leyes de escalado predecibles sin evidencia de saturación. El \textit{flow matching} \citep{lipman2022flowmatching,liu2022rectified} define trayectorias lineales entre el ruido y los datos:

\begin{equation}
\mathcal{L}_{\text{FM}} = \E_{t, \x_0, \eps}\!\left[\left\|v_\theta((1-t)\x_0 + t\eps, t) - (\eps - \x_0)\right\|_2^2\right]
\end{equation}

\section{Estado del arte en análisis del espacio latente}
\label{sec:sota}

\subsection{El espacio $\mathcal{W}$ de StyleGAN y la hipótesis de linealidad}

El análisis estructural del espacio latente encontró en StyleGAN \citep{karras2019stylegan} su objeto de estudio más fructífero. La hipótesis de linealidad, que habla sobre cómo los conceptos semánticos corresponden a regiones linealmente separables, ha demostrado ser empíricamente robusta en $\mathcal{W}$. InterFaceGAN \citep{shen2020interfacegan} entrenó una SVM lineal para encontrar el hiperplano que separa imágenes con y sin un atributo dado:

\begin{equation}
\langle \mathbf{n}, \mathbf{w} \rangle + b = 0
\end{equation}

La edición se realiza como $\mathbf{w}' = \mathbf{w} + \alpha\mathbf{n}$. La contribución más relevante es que estas direcciones son transferibles: vectores encontrados sobre un subconjunto de latentes se aplican con alta calidad a latentes arbitrarios, confirmando que corresponden a propiedades globales de la distribución aprendida.

GANSpace \citep{harkonen2020ganspace} demostró que el PCA sobre vectores $\mathbf{w}$ produce componentes interpretables de forma no supervisada. SeFa \citep{shen2021sefa} mostró que los vectores propios de la primera capa del generador son suficientes para recuperar estas direcciones sin necesidad de datos. \citet{wu2021stylespace} propusieron el \textit{StyleSpace} $\mathcal{S}$ formado por los coeficientes de modulación canal a canal, con un disentanglement superior al de $\mathcal{W}$.

\subsection{El espacio latente en los modelos de difusión latente}

La inversión DDIM \citep{song2021ddim} proporciona la base técnica para recuperar el latente que genera una imagen dada; \citet{mokady2023nulltext} refinaron este procedimiento con la \textit{inversión de texto nulo} para garantizar reconstrucciones exactas. \citet{kwon2023diffusion} identificaron el \textit{h-space} (activaciones del bottleneck del UNet) como un espacio con linealidad semántica análoga al $\mathcal{W}$ de StyleGAN:

\begin{equation}
\mathbf{h}_t = \text{UNet}_\text{enc}(\z_t, t, \tau_\theta(y))
\end{equation}

La operación de edición $\mathbf{h}'_t = \mathbf{h}_t + \alpha\Delta\mathbf{h}$ produce cambios semánticos coherentes que transfieren a imágenes no vistas. Sin embargo, el trabajo se limita a un único modelo sin estudiar la dependencia con la arquitectura. \citet{hertz2022p2p} demostraron que los mapas de atención cruzada codifican la disposición espacial de los conceptos, habilitando Prompt-to-Prompt. \citet{tang2023emergent} identificaron que las activaciones del UNet en pasos de tiempo intermedios superan a DINO y CLIP para correspondencia semántica densa.

\subsection{Dimensionalidad intrínseca de los espacios latentes}

\citet{pope2021intrinsic} estudiaron la dimensionalidad intrínseca de los espacios latentes usando el estimador de dos vecinos más cercanos (TwoNN, \citealp{facco2017twonn}), encontrando que el espacio $\mathcal{W}$ de StyleGAN para rostros tiene una dimensionalidad intrínseca de aproximadamente 40--60 dimensiones efectivas pese a sus 512 nominales. Esta propiedad de baja dimensionalidad es una hipótesis verificable también para los LDMs, cuya dimensión nominal oscila entre 16.384 y 65.536 elementos.

\subsection{Fine-tuning y personalización}

La inversión textual \citep{gal2022textual} optimiza el embedding de un token virtual sin modificar los pesos del modelo. DreamBooth \citep{ruiz2023dreambooth} fine-tunea el UNet completo con pérdida de preservación de clases para evitar el olvido catastrófico \citep{mccloskey1989catastrophic,kirkpatrick2017ewc}. LoRA \citep{hu2022lora} factoriza el desplazamiento de pesos en bajo rango:

\begin{equation}
\mathbf{W} = \mathbf{W}_0 + \mathbf{B}\mathbf{A},
\quad \mathbf{B} \in \R^{d \times r},\; \mathbf{A} \in \R^{r \times k},\; r \ll \min(d, k)
\end{equation}

El análisis de cómo el fine-tuning modifica la estructura interna del espacio latente es el aspecto de mayor novedad de este trabajo respecto a la literatura existente.

\subsection{Brechas identificadas y posicionamiento}

La revisión permite identificar cuatro brechas principales. Primera: la práctica totalidad de los estudios de análisis del espacio latente en LDMs se realizan sobre SD~1.x, sin comparación sistemática entre arquitecturas. Segunda: la literatura se centra en el h-space y los mapas de atención, no en el espacio $\z_0$ producido por el VAE, más accesible para aplicaciones de edición basadas en inversión DDIM. Tercera: los métodos de personalización se han evaluado en términos de calidad visual, pero no en términos de las propiedades estructurales del espacio latente resultante. Cuarta: no existe un protocolo replicable para la evaluación comparativa del espacio latente de LDMs aplicable uniformemente a distintas arquitecturas.

\section{Objetivos}
\label{sec:objetivos}

El objetivo general es caracterizar y comparar cuantitativamente la estructura del espacio latente $\z_0$ de SD~1.5, SD~2.1 y SDXL, y analizar cómo dicha estructura se modifica tras el fine-tuning sobre un dominio visual específico. Los objetivos específicos son: (1) cuantificar la relación entre distancia latente y similitud perceptual y caracterizar la dimensionalidad intrínseca; (2) identificar y cuantificar direcciones semánticas para atributos visuales controlados y evaluar el grado de disentanglement; (3) medir la deriva geométrica y semántica del espacio latente tras LoRA sobre WikiArt y comparar la respuesta de los tres modelos; y (4) identificar qué propiedades del espacio latente son invariantes a la arquitectura y cuáles dependen del codificador de texto, la capacidad del UNet o la resolución nativa.

% ===========================================================
%  METODOLOGÍA
% ===========================================================
\chapter{Metodología}
\label{ch:metodos}

\section{Entorno experimental}
\label{sec:entorno}

\subsection{Plataforma de cómputo}

Los experimentos se ejecutaron sobre Google Colaboratory \citep{google2017colab}, utilizando la GPU NVIDIA Tesla T4 (15 GB VRAM, arquitectura Turing). La decisión de emplear computación en la nube responde a las limitaciones del equipo local (NVIDIA GeForce GTX~1050~Max-Q, 3~GB~VRAM), insuficiente para ejecutar los modelos de difusión en condiciones reproducibles. La persistencia de datos entre sesiones se gestiona mediante Google Drive. El equipo local (Intel Core i5-1035G1, 7,5~GB~RAM, Ubuntu~22.04) se utiliza exclusivamente para el desarrollo del código.

\subsection{Entorno de software y reproducibilidad}

La Tabla~\ref{tab:software} recoge las versiones exactas de las dependencias principales. La reproducibilidad se garantiza mediante especificación explícita de versiones en requirements\_colab.txt, un cuaderno de entrada estandarizado que instala las dependencias al inicio de cada sesión, control de versiones con Git, y semillas aleatorias fijas en PyTorch, NumPy y Python.

\begin{table}[htb]
\centering
\caption{Versiones exactas de las dependencias del entorno experimental.}
\label{tab:software}
\small
\begin{tabular}{lll}
\toprule
\textbf{Biblioteca} & \textbf{Versión} & \textbf{Propósito} \\
\midrule
Python       & 3.11.15       & Lenguaje base \\
PyTorch      & 2.3.1+cu121   & Framework de aprendizaje profundo \\
Diffusers    & 0.27.2        & Pipelines de modelos de difusión \\
Transformers & 4.40.2        & Codificadores de texto \\
Accelerate   & 0.30.1        & Entrenamiento multi-dispositivo \\
PEFT         & 0.11.1        & Fine-tuning eficiente (LoRA) \\
xformers     & 0.0.27        & Atención eficiente en memoria \\
open-clip-torch & 2.24.0    & Modelo CLIP de código abierto \\
scikit-learn & 1.8.0         & Análisis estadístico y PCA \\
\bottomrule
\end{tabular}
\end{table}

\subsection{Métricas de evaluación}

Se utilizan dos métricas de similitud de imagen. El MSE mide la diferencia cuadrática media píxel a píxel y sirve como referencia global de disimilitud. El SSIM global \citep{wang2004ssim} evalúa la similitud estructural entre las estadísticas de las imágenes y constituye la métrica perceptual principal del análisis de correlación latente--imagen.

\section{Modelos bajo estudio}
\label{sec:modelos}

\subsection{Marco teórico: modelos de difusión latente}

Los modelos de difusión latente \citep{rombach2022ldm} operan en dos etapas: compresión perceptual mediante un VAE y difusión en el espacio latente comprimido. El proceso de difusión en el espacio latente sigue:

\begin{equation}
q(\z_t \mid \z_0) = \mathcal{N}\!\left(\sqrt{\bar{\alpha}_t}\,\z_0,\; (1 - \bar{\alpha}_t)\,\mathbf{I}\right)
\end{equation}

El modelo de ruido $\eps_\theta$ se entrena minimizando:

\begin{equation}
\mathcal{L} = \E_{\z_0,\, \eps \sim \mathcal{N}(0,\mathbf{I}),\, t}\!
\left[\norm{\eps - \eps_\theta(\z_t, t, \tau_\theta(y))}^2\right]
\end{equation}

donde $\tau_\theta(y)$ es la representación del prompt producida por el codificador de texto.

\subsection{Criterios de selección}

La selección se guió por tres criterios. Diversidad arquitectónica: los modelos difieren en el codificador de texto y en la capacidad del UNet, las dos variables que se desean estudiar. VAE compartido: los tres modelos utilizan la misma arquitectura de autoencoder variacional con factor de compresión $f = 8$ y $C = 4$ canales latentes, garantizando que las diferencias observadas en el espacio $\z_0$ son atribuibles al proceso de difusión y al condicionamiento textual. Viabilidad computacional dentro de los 15~GB~VRAM de la GPU~T4.

\subsection{Stable Diffusion 1.5}

SD~1.5 \citep{rombach2022ldm} fue entrenado por RunwayML sobre LAION-Aesthetics v2 5+. El codificador de texto es CLIP ViT-L/14 \citep{radford2021clip}, que proyecta el prompt en 77 tokens de 768 dimensiones. El UNet tiene aproximadamente 860 millones de parámetros y la resolución nativa es $512 \times 512$ con latentes de forma $4 \times 64 \times 64$.

\subsection{Stable Diffusion 2.1}

SD~2.1 \citep{stabilityai2022sd2} introduce dos diferencias respecto a SD~1.5. El codificador OpenCLIP ViT-H/14 \citep{ilharco2021openclip} produce embeddings de 1.024 dimensiones frente a las 768 de CLIP ViT-L/14. La resolución nativa es $768 \times 768$ con latentes de forma $4 \times 96 \times 96$.

\subsection{Stable Diffusion XL}

SDXL \citep{podell2023sdxl} introduce tres cambios. El doble codificador de texto CLIP ViT-L/14 más OpenCLIP ViT-bigG/14 (1.300 millones de parámetros) produce un vector de condicionamiento de 2.048 dimensiones de atención cruzada más 1.280 de embedding global. El UNet de 2.600 millones de parámetros es $\sim 3\times$ mayor que el de los modelos anteriores. La resolución nativa de $1024 \times 1024$ produce latentes de forma $4 \times 128 \times 128$.

\begin{table}[htb]
\centering
\caption{Comparativa arquitectónica de los modelos bajo estudio.}
\label{tab:arquitecturas}
\small
\begin{tabular}{lccc}
\toprule
\textbf{Propiedad} & \textbf{SD 1.5} & \textbf{SD 2.1} & \textbf{SDXL} \\
\midrule
Codificador de texto & CLIP ViT-L/14 & OpenCLIP ViT-H/14 & CLIP ViT-L + OpenCLIP ViT-bigG \\
Dim.\ embedding texto & 768 & 1.024 & 2.048 + 1.280 \\
Parámetros UNet (aprox.) & 860 M & 870 M & 2.600 M \\
Resolución nativa & $512\times512$ & $768\times768$ & $1024\times1024$ \\
Factor de compresión ($f$) & 8 & 8 & 8 \\
Canales latentes ($C$) & 4 & 4 & 4 \\
Forma del latente & $4\times64\times64$ & $4\times96\times96$ & $4\times128\times128$ \\
Dimensión latente total & 16.384 & 36.864 & 65.536 \\
\bottomrule
\end{tabular}
\end{table}

\subsection{Interfaz unificada y extracción de $\z_0$}

Se implementó la clase ModelWrapper que encapsula el pipeline de Diffusers y expone tres operaciones estables: codificación de imagen a latente, decodificación de latente a imagen, y generación con captura del latente final. El latente $\z_0$ en modo codificación se extrae como la media de la distribución posterior multiplicada por el factor de escalado del VAE, garantizando determinismo. En modo generación se usa un hook temporal sobre vae.decode para capturar $\z_0$ antes de la decodificación.

\section{Corpus de vectores latentes}
\label{sec:corpus}

\subsection{Diseño}

Se generaron 150 vectores latentes por modelo, organizados en dos categorías latentes base (60 por modelo, 20 prompts $\times$ 3 semillas) para el análisis geométrico global y la base PCA del análisis semántico; y latentes de atributo (90 por modelo, 15 grupos $\times$ 2 valores $\times$ 3 semillas) para construir las direcciones semánticas. Los mismos prompts, semillas e hiperparámetros se aplican a los tres modelos.

\subsection{Diseño del banco de prompts}

Los 20 prompts base y los 30 prompts de atributo se diseñaron con el asistente de inteligencia artificial Claude \citep{anthropic2024claude}, empleándolo para generar descripciones diversas, equilibradas y libres de correlación temática inadvertida. Los prompts base cubren una amplia diversidad de dominios visuales (naturaleza muerta, fauna, arquitectura, gastronomía, paisaje, espacio urbano, tecnología), escalas compositivas y condiciones de luz. La lista completa se recoge en el Apéndice~\ref{app:prompts}.

Los 15 grupos de atributo se construyeron como pares controlados: ambos prompts del mismo grupo describen la misma escena modificando exclusivamente el atributo objetivo. Los tres tipos de atributo son color (temperatura cálida frente a fría, 5 grupos), iluminación (alta clave frente a baja clave, 5 grupos) y estilo (fotografía fotorrealista frente a imagen artística, 5 grupos). La dirección semántica de un grupo $g$ se define como:

\begin{equation}
\mathbf{d}_g = \bar{\z}(g, v^+) - \bar{\z}(g, v^-)
\label{eq:dir_semantica}
\end{equation}

donde $\bar{\z}(g, v)$ es la media sobre las tres semillas, eliminando la variabilidad estocástica para aislar el efecto del atributo.

\subsection{Hiperparámetros de generación}

\begin{table}[htb]
\centering
\caption{Hiperparámetros de generación del corpus.}
\label{tab:hiper_generacion}
\small
\begin{tabular}{lccc}
\toprule
\textbf{Parámetro} & \textbf{SD 1.5} & \textbf{SD 2.1} & \textbf{SDXL} \\
\midrule
Sampler & DDIM & DDIM & DDIM \\
Pasos de denoising & 50 & 50 & 30 \\
Guidance scale $w$ & 7,5 & 7,5 & 7,5 \\
Semillas & 0, 1, 2 & 0, 1, 2 & 0, 1, 2 \\
Resolución de salida & $512^2$ & $512^2$ & $512^2$ \\
Precisión & fp16 & fp16 & fp16 (VAE float32) \\
\bottomrule
\end{tabular}
\end{table}

\section{Análisis geométrico}
\label{sec:metodos:geometrico}

\subsection{Separabilidad intra/inter-prompt}

Para cada modelo se construye la matriz de distancias euclídeas $\mathbf{D} \in \R^{N \times N}$ sobre los $N = 60$ latentes base. El ratio de separabilidad $\text{sep} = \mu_{\text{inter}} / \mu_{\text{intra}}$ mide si los latentes del mismo prompt están más agrupados que los de prompts distintos. La significación estadística se evalúa con el test de Kolmogorov--Smirnov bilateral.

\subsection{Dimensionalidad efectiva}

Los valores propios de la covarianza se obtienen mediante SVD thin de la matriz de datos centrada. El participation ratio:

\begin{equation}
\text{PR} = \frac{\left(\sum_k \lambda_k\right)^2}{\sum_k \lambda_k^2}
\end{equation}

mide el número efectivo de dimensiones activas de forma continua y sin necesidad de elegir un umbral.

\subsection{Correlación distancia latente--similitud imagen}

Se miden la correlación de Pearson $r$ y de Spearman $\rho$ entre la distancia latente L2 y las métricas de disimilitud MSE y $-\text{SSIM}$.

\section{Análisis semántico}
\label{sec:metodos:semantico}

\subsection{Análisis de componentes principales}

El PCA se aplica sobre los 60 latentes base aplanados a vectores de dimensión $D$. Las métricas extraídas son la varianza explicada por componente y el número de componentes para el 80\% y 90\% de varianza acumulada.

\subsection{Consistencia y ortogonalidad de las direcciones semánticas}

La consistencia intra-tipo mide si la dirección de un atributo es universal (escena-independiente):

\begin{equation}
\text{consistencia}(t) = \frac{1}{\binom{|G_t|}{2}} \sum_{g \neq g' \in G_t} \cos\!\left(\hat{\mathbf{d}}_g, \hat{\mathbf{d}}_{g'}\right)
\end{equation}

El entrelazamiento inter-tipo mide la ausencia de disentanglement:

\begin{equation}
\text{entrelazamiento}(t_1, t_2) = \frac{1}{|G_{t_1}| \cdot |G_{t_2}|}
\sum_{g,\, g'} \left|\cos\!\left(\hat{\mathbf{d}}_g, \hat{\mathbf{d}}_{g'}\right)\right|
\end{equation}

La fracción de energía semántica en los $K_0 = 10$ primeros componentes PCA cuantifica si el atributo se codifica en el subespacio de alta varianza:

\begin{equation}
f_g = \frac{\sum_{k=1}^{K_0} (\mathbf{V}_k^\top \hat{\mathbf{d}}_g)^2}{\sum_{k=1}^{K} (\mathbf{V}_k^\top \hat{\mathbf{d}}_g)^2}
\end{equation}

\section{Fine-tuning y análisis de deriva latente}
\label{sec:metodos:finetune}

\subsection{Diseño experimental}

El análisis de fine-tuning cierra el ciclo del trabajo: los espacios latentes caracterizados en las secciones anteriores se someten a adaptación de dominio y se vuelven a caracterizar con los mismos instrumentos. Los tres modelos se adaptan sobre el mismo dataset, con los mismos hiperparámetros y el mismo número de pasos de actualización.

\subsection{Selección del dominio: WikiArt}

Se selecciona el dominio de pintura clásica utilizando WikiArt \citep{saleh2015wikiart} por cuatro razones. Primero, ataca directamente el atributo con consistencia intra-tipo más baja en los tres modelos (estilo: 0,028 en SD~1.5, 0,010 en SD~2.1, 0,209 en SDXL), permitiendo verificar si la exposición intensiva a un dominio artístico induce la aparición de una dirección de estilo más universal. Segundo, la distribución de pinturas artísticas es cualitativamente diferente de la fotografía dominante en LAION-5B \citep{schuhmann2022laion}, maximizando la señal detectable en las métricas de deriva. Tercero, dado que SDXL ya tiene consistencia de estilo un orden de magnitud mayor que SD~1.5 y SD~2.1, existe una hipótesis diferencial comprobable: \textit{SDXL absorberá el conocimiento artístico de WikiArt de forma más eficiente y producirá mayor amplificación de la consistencia semántica de estilo.} Cuarto, WikiArt es un dataset de referencia en la literatura de síntesis de estilos artísticos \citep{gatys2016style}.

\subsection{Dataset WikiArt}

Se selecciona un subconjunto balanceado de 500 imágenes: 50 por cada uno de los diez movimientos artísticos representados (impresionismo, realismo, romanticismo, expresionismo, post-impresionismo, barroco, cubismo, arte abstracto, simbolismo y Art Nouveau). Las captions se generan a partir de los metadatos con la plantilla \textit{``a \{style\} painting\{, \{genre\},\} by \{artist\}''}.
\begin{figure}[htb]
\centering
\includegraphics[width=\textwidth]{wikiart_dataset_preview}
\caption{Muestra representativa del subconjunto de WikiArt utilizado para el fine-tuning (500 imágenes, 50 por movimiento artístico). Los diez movimientos incluidos son: impresionismo, realismo, romanticismo, expresionismo, post-impresionismo, barroco, cubismo, arte abstracto, simbolismo y Art~Nouveau. La distribución visual es cualitativamente opuesta a la fotografía naturalista de LAION-5B, maximizando la señal detectable en las métricas de deriva latente.}
\label{fig:wikiart_preview}
\end{figure}

\subsection{Configuración LoRA}

LoRA se aplica exclusivamente al UNet; el codificador de texto y el VAE se mantienen congelados. La congelación del codificador de texto asegura que los cambios en los latentes sean atribuibles a la reorganización del proceso de denoising. La congelación del VAE garantiza que $\z_0^\text{pre}$ y $\z_0^\text{ft}$ sean comparables en el mismo espacio vectorial.

\begin{table}[htb]
\centering
\caption{Configuración LoRA y hiperparámetros de entrenamiento.}
\label{tab:lora}
\small
\begin{tabular}{lccc}
\toprule
\textbf{Parámetro} & \textbf{SD 1.5} & \textbf{SD 2.1} & \textbf{SDXL} \\
\midrule
Rango ($r$) & 8 & 8 & 8 \\
Alpha ($\alpha$) & 8 & 8 & 8 \\
Dropout & 0,1 & 0,1 & 0,1 \\
Módulos objetivo & \multicolumn{3}{c}{\texttt{to\_q, to\_k, to\_v, to\_out.0}} \\
Parámetros entrenables & $\approx 3{,}4$ M & $\approx 3{,}5$ M & $\approx 11{,}2$ M \\
Fracción del total & 0,4\% & 0,4\% & 0,4\% \\
Pasos de entrenamiento & 1.000 & 1.000 & 800 \\
Batch efectivo & 8 & 8 & 4 \\
Learning rate & $10^{-4}$ & $10^{-4}$ & $10^{-4}$ \\
LR scheduler & cosine & cosine & cosine \\
Optimizador & AdamW & AdamW & AdamW \\
\bottomrule
\end{tabular}
\end{table}

\subsection{Métricas de deriva latente}

Tras el fine-tuning se regenera el corpus completo con el modelo adaptado. Las métricas de deriva son: $\Delta\text{sep}$, $\Delta\text{PR}$ y drift del centroide para la deriva geométrica; alineación de Grassmann entre los $K_0 = 10$ subespacios principales para la rotación PCA; y $\Delta\norm{\mathbf{d}_g}$, $\Delta\text{consistencia}$, $\Delta\text{entrelazamiento}$ y $\Delta f_g$ para la deriva semántica.

% ===========================================================
%  RESULTADOS
% ===========================================================
\chapter{Resultados}
\label{ch:resultados}

El capítulo organiza los resultados en cuatro bloques que siguen el orden experimental del trabajo.
La sección~\ref{sec:res:modelos} reporta la verificación de los tres modelos y las estadísticas
descriptivas del espacio latente $\z_0$. La sección~\ref{sec:res:corpus} documenta el corpus
generado. Las secciones~\ref{sec:res:geometrico} y \ref{sec:res:semantico} presentan el análisis
geométrico y semántico del pretrain. La sección~\ref{sec:res:finetune} cierra el capítulo con los
resultados del fine-tuning sobre WikiArt y el análisis de deriva latente.

% -----------------------------------------------------------
\section{Verificación de modelos y análisis del espacio latente $\z_0$}
\label{sec:res:modelos}
% -----------------------------------------------------------

\subsection{Carga y estadísticas descriptivas}

Los tres modelos se cargaron correctamente en la GPU T4 en precisión fp16.
La excepción es el VAE de SDXL, que requiere float32 para evitar desbordamiento
numérico durante la decodificación, motivado por los valores de activación de mayor magnitud
del codificador de mayor resolución~\citep{podell2023sdxl}.
La Tabla~\ref{tab:estadisticas_z0} recoge las estadísticas descriptivas del latente $\z_0$
para el mismo prompt y semilla en los tres modelos.

\begin{table}[htb]
\centering
\caption{Estadísticas descriptivas del espacio latente $\z_0$ (mismo prompt y semilla).
$\bar{z}$: media; $\sigma_z$: desviación típica; $\norm{\z}$: norma L2; MAE~VAE: error
absoluto medio del ciclo codificación\textrightarrow{}decodificación en escala de píxel $[0,255]$.}
\label{tab:estadisticas_z0}
\small
\begin{tabular}{lrrrrr}
\toprule
\textbf{Modelo} & \textbf{Dim.} & $\bar{z}$ & $\sigma_z$ & $\norm{\z}$ & \textbf{MAE VAE} \\
\midrule
SD 1.5 & 16.384 & $-0{,}083$ & $0{,}856$ & $110{,}0$ & $4{,}00$ \\
SD 2.1 & 36.864 & $-0{,}162$ & $0{,}669$ & $132{,}2$ & $2{,}64$ \\
SDXL   & 65.536 & $-0{,}188$ & $0{,}595$ & $159{,}7$ & $1{,}80$ \\
\bottomrule
\end{tabular}
\end{table}

La norma L2 crece monótonamente a pesar de que la desviación típica decrece, coherente con
la expresión $\E[\norm{\z}] \approx \sqrt{N(\mu^2 + \sigma^2)}$ bajo distribución aproximadamente
gaussiana, confirmada visualmente por los histogramas de cada canal.
La mejora progresiva del MAE del ciclo encode\textrightarrow{}decode ($4{,}00 \to 2{,}64 \to 1{,}80$)
indica que los VAE de los modelos más recientes poseen mayor capacidad; el error del VAE de SDXL
(0,71\% relativo respecto al rango $[0,255]$) es subperceptual y valida la idoneidad del espacio
$\z_0$ como representación de análisis.

\subsection{Estructura espacial de los canales latentes}

\begin{figure}[htb]
\centering
\includegraphics[width=\textwidth]{data/fase1/fase1_latent_channels}
\caption{Visualización de los cuatro canales del tensor $\z_0$ para SD~1.5 (izquierda), SD~2.1
(centro) y SDXL (derecha) con el mismo prompt y semilla. El rango de valores aumenta de
$[-6,+4]$ en SD~1.5 a $[-10,+7{,}5]$ en SD~2.1.}
\label{fig:canales}
\end{figure}

La inspección visual (Figura~\ref{fig:canales}) revela diferencias cualitativas entre modelos.
En SD~1.5, los cuatro canales muestran la silueta del objeto principal con complementariedad
sistemática entre pares. En SD~2.1, la estructura semántica se conserva con mayor rango de valores,
reflejo del condicionamiento más expresivo del codificador de mayor dimensión. En SDXL, destaca
la presencia de bandas verticales periódicas que corresponden al patrón estructural del fondo de la
imagen generada: el VAE de SDXL codifica los patrones espaciales del fondo con una fidelidad
comparable al objeto principal, evidenciando una representación más distribuida y coherente con su
menor MAE.

\subsection{Verificación del ciclo de reconstrucción}

\begin{figure}[htb]
\centering
\begin{subfigure}[b]{0.32\textwidth}
    \includegraphics[width=\textwidth]{data/fase1/sd15/roundtrip}
    \caption*{SD~1.5 --- MAE $= 4{,}00$}
\end{subfigure}
\hfill
\begin{subfigure}[b]{0.32\textwidth}
    \includegraphics[width=\textwidth]{data/fase1/sd21/roundtrip}
    \caption*{SD~2.1 --- MAE $= 2{,}64$}
\end{subfigure}
\hfill
\begin{subfigure}[b]{0.32\textwidth}
    \includegraphics[width=\textwidth]{data/fase1/sdxl/roundtrip}
    \caption*{SDXL --- MAE $= 1{,}80$}
\end{subfigure}
\caption{Comparativa del ciclo codificación\textrightarrow{}decodificación (roundtrip) para los
tres modelos. La imagen generada se codifica a $\z_0$ y se decodifica de nuevo; la diferencia
absoluta media confirma la idoneidad del VAE como representación fiel del contenido visual.}
\label{fig:roundtrip}
\end{figure}

La Figura~\ref{fig:roundtrip} muestra el ciclo de reconstrucción completo para cada modelo.
Las diferencias visibles son mínimas en SDXL y apreciables solo en texturas de alta frecuencia
en SD~1.5. La reducción del error en un factor 2,2 entre SD~1.5 y SDXL confirma que la
mayor dimensión latente de SDXL ($4\times128\times128$ frente a $4\times64\times64$) preserva
más información de alta frecuencia, con independencia del proceso de difusión.

% -----------------------------------------------------------
\section{Corpus de vectores latentes}
\label{sec:res:corpus}
% -----------------------------------------------------------

Se generaron 450 vectores latentes en total (150 por modelo), almacenados como tensores
float32 en formato $[N, 1, C, H, W]$ con los mismos prompts, semillas e hiperparámetros
aplicados a los tres modelos de forma estrictamente sincronizada.
La verificación de reproducibilidad confirmó que la regeneración de un subconjunto de muestras
con los mismos hiperparámetros produce vectores idénticos bit a bit, validando la determinismo
del muestreo DDIM~\citep{song2021ddim} con semillas fijas.

% -----------------------------------------------------------
\section{Análisis geométrico del espacio latente}
\label{sec:res:geometrico}
% -----------------------------------------------------------

\subsection{Separabilidad intra/inter-prompt}
\label{subsec:res:separabilidad}

\begin{table}[htb]
\centering
\caption{Estadísticas de distancias L2 intra e inter-prompt (60 latentes base por modelo).
$\mu_\text{intra}$: distancia media entre latentes del mismo prompt con semillas distintas;
$\mu_\text{inter}$: distancia media entre latentes de prompts distintos;
sep $= \mu_\text{inter} / \mu_\text{intra}$; KS: estadístico de Kolmogorov--Smirnov bilateral;
$p$: nivel de significación.}
\label{tab:separabilidad}
\small
\begin{tabular}{lrrrrrrr}
\toprule
\textbf{Modelo} & $\mu_\text{intra}$ & $\sigma_\text{intra}$ & $\mu_\text{inter}$ &
$\sigma_\text{inter}$ & \textbf{sep} & \textbf{KS} & \textbf{$p$} \\
\midrule
SD 1.5 & 151,8 & 16,4 & 148,1 & 22,4 & 0,976 & 0,142 & 0,173 \\
SD 2.1 & 195,8 & 22,4 & 189,1 & 30,7 & 0,966 & 0,187 & 0,031 \\
SDXL   & 230,5 & 35,0 & 257,6 & 42,9 & \textbf{1,118} & 0,304 & $<{0{,}001}$ \\
\bottomrule
\end{tabular}
\end{table}

\begin{figure}[htbp]
\centering
\begin{subfigure}[b]{\textwidth}
    \includegraphics[width=\textwidth]{data/fase3/results/geo_sd15_distancias}
    \caption*{SD~1.5}
\end{subfigure}\\[0.4em]
\begin{subfigure}[b]{\textwidth}
    \includegraphics[width=\textwidth]{data/fase3/results/geo_sd21_distancias}
    \caption*{SD~2.1}
\end{subfigure}\\[0.4em]
\begin{subfigure}[b]{\textwidth}
    \includegraphics[width=\textwidth]{data/fase3/results/geo_sdxl_distancias}
    \caption*{SDXL}
\end{subfigure}
\caption{Distribuciones de distancias L2 intra-prompt e inter-prompt. En SD~1.5 y SD~2.1 los
histogramas se superponen casi completamente (sep $< 1$); en SDXL están claramente separados
(sep $> 1$, $p < 0{,}001$).}
\label{fig:distancias}
\end{figure}

SD~1.5 y SD~2.1 presentan sep $< 1$ (0,976 y 0,966), indicando que los latentes del mismo
prompt con distintas semillas están en promedio \emph{más alejados} que latentes de prompts
distintos (Tabla~\ref{tab:separabilidad}, Figura~\ref{fig:distancias}). El test KS no rechaza
la igualdad de distribuciones para SD~1.5 ($p = 0{,}173$) y solo marginalmente para SD~2.1
($p = 0{,}031$). En estos modelos, la variabilidad estocástica del proceso de denoising
desplaza el latente tanto como un cambio de contenido semántico completo: la relación
señal-ruido es desfavorable para las técnicas de edición basadas en aritmética de vectores
en $\z_0$.

SDXL, en cambio, presenta sep $= 1{,}118$ con $p < 0{,}001$: los latentes del mismo prompt
están sistemáticamente más próximos que los de prompts distintos. El contenido del prompt
es el predictor dominante de la posición del latente, con la variabilidad estocástica
subordinada al significado semántico.

\subsection{Dimensionalidad efectiva e inspección PCA}
\label{subsec:res:dimensionalidad}

Los valores propios de la covarianza se obtienen mediante SVD thin de la matriz de datos
centrada sobre los 60 latentes base aplanados a vectores de dimensión $D$ (16.384, 36.864
o 65.536 elementos según el modelo). La Tabla~\ref{tab:dimensionalidad} reúne el participation
ratio (PR), los umbrales de varianza acumulada y las fracciones explicadas por las dos primeras
componentes, evitando así la repetición de información entre tablas.

\begin{table}[htb]
\centering
\caption{Dimensionalidad efectiva e inspección de los primeros componentes PCA (60 latentes
base por modelo). PR: participation ratio; dims@80/90/95\%: número de componentes para
alcanzar el umbral de varianza acumulada indicado; PC1, PC2: fracción de varianza explicada
por el primero y segundo componente respectivamente.}
\label{tab:dimensionalidad}
\small
\begin{tabular}{lrrrrrrr}
\toprule
\textbf{Modelo} & \textbf{Dim.\ nom.} & \textbf{PR} & \textbf{dims@80\%} &
\textbf{dims@90\%} & \textbf{dims@95\%} & \textbf{PC1 (\%)} & \textbf{PC2 (\%)} \\
\midrule
SD 1.5 & 16.384 & 12,35 & 25 & 38 & 46 & 19,2 & 17,7 \\
SD 2.1 & 36.864 & 12,33 & 27 & 39 & 47 & 20,1 & 17,1 \\
SDXL   & 65.536 & 11,56 & 22 & 35 & 45 & 19,1 & 14,6 \\
\bottomrule
\end{tabular}
\end{table}

\begin{figure}[htb]
\centering
\includegraphics[width=0.85\textwidth]{data/fase3/results/snapshot_cross_model_pca}
\caption{Varianza acumulada en función del número de componentes principales para los tres
modelos. Las curvas son prácticamente indistinguibles hasta el componente 10, donde SDXL
diverge ligeramente hacia una mayor concentración.}
\label{fig:pca_varianza}
\end{figure}

El participation ratio PR $\approx 12$ es prácticamente idéntico en los tres modelos pese
a que sus dimensiones nominales difieren en un factor de cuatro
(Tabla~\ref{tab:dimensionalidad}, Figura~\ref{fig:pca_varianza}).
El número de componentes para capturar el 80\% de la varianza (22--27) representa
menos del 0,04\% de la dimensión nominal de SDXL, confirmando que la baja dimensionalidad
intrínseca es una propiedad arquitectónica compartida, no vinculada al tamaño del espacio.

La diferencia más notable reside en la tercera componente de SDXL, que explica el 13,3\%
de la varianza (acumulado del 47,1\% con tres componentes), frente al 6,1\% y 5,3\% en los
modelos anteriores. Esa tercera componente coincide con la dirección semántica de
iluminación (véase la sección~\ref{sec:res:semantico}).

\subsection{Correlación entre distancia latente y similitud de imagen}
\label{subsec:res:correlacion}

Se muestrearon aleatoriamente 500 pares de los $\binom{60}{2} = 1.770$ posibles (semilla fija 42)
y se calcularon las correlaciones de Pearson ($r$) y Spearman ($\rho$) entre la distancia L2
latente y las métricas de disimilitud MSE y $-\text{SSIM}$~\citep{wang2004ssim}.

\begin{table}[htb]
\centering
\caption{Correlación entre distancia latente L2 y métricas de disimilitud de imagen
(500 pares por modelo, semilla 42).}
\label{tab:correlacion}
\small
\begin{tabular}{lrrrr}
\toprule
\textbf{Modelo} & $r(\text{L2, MSE})$ & $\rho(\text{L2, MSE})$ &
$r(\text{L2}, -\text{SSIM})$ & \textbf{$p$-valor} \\
\midrule
SD 1.5 & 0,740 & 0,763 & 0,796 & $\ll 10^{-80}$ \\
SD 2.1 & 0,718 & 0,753 & 0,761 & $\ll 10^{-80}$ \\
SDXL   & 0,680 & 0,698 & 0,691 & $\ll 10^{-68}$ \\
\bottomrule
\end{tabular}
\end{table}

La correlación es positiva, fuerte y altamente significativa en los tres modelos
(Tabla~\ref{tab:correlacion}). El orden $r_{\text{sd15}} > r_{\text{sd21}} > r_{\text{sdxl}}$
sugiere que el VAE de menor resolución de SD~1.5 actúa como filtro de información de alta
frecuencia más agresivo, produciendo latentes cuya distancia L2 está más alineada con
métricas globales como el MSE. El mayor MAE del ciclo roundtrip de SD~1.5 y la mayor
correlación con MSE son dos caras del mismo fenómeno: el VAE de SD~1.5 descarta más
información de alta frecuencia, haciendo que las distancias latentes reflejen más fielmente
las diferencias semánticas globales capturadas por el MSE.

% -----------------------------------------------------------
\section{Análisis semántico del espacio latente}
\label{sec:res:semantico}
% -----------------------------------------------------------

\subsection{Discriminabilidad de las direcciones semánticas}
\label{subsec:res:discriminabilidad}

La dirección semántica de cada grupo de atributo $g$ se define como la diferencia de medias
(Ec.~\ref{eq:dir_semantica}, §3.4). La Tabla~\ref{tab:normas_dir} recoge las normas medias
por tipo de atributo y modelo.

\begin{table}[htb]
\centering
\caption{Norma media ($\pm$ desviación típica) de las direcciones semánticas
$\mathbf{d}_g = \bar{\z}(g,v^+) - \bar{\z}(g,v^-)$ por tipo de atributo y modelo.}
\label{tab:normas_dir}
\small
\begin{tabular}{lrrr}
\toprule
\textbf{Tipo de atributo} & \textbf{SD 1.5} & \textbf{SD 2.1} & \textbf{SDXL} \\
\midrule
Color (cálido/frío)            & $65{,}4 \pm 4{,}1$  & $90{,}4 \pm 10{,}9$ & $183{,}4 \pm 18{,}8$ \\
Iluminación (brillante/oscuro) & $55{,}8 \pm 11{,}4$ & $76{,}1 \pm 13{,}5$ & $149{,}8 \pm 29{,}1$ \\
Estilo (foto/artístico)        & $51{,}6 \pm 15{,}0$ & $64{,}8 \pm 12{,}3$ & $106{,}5 \pm 32{,}7$ \\
\bottomrule
\end{tabular}
\end{table}

\begin{figure}[htbp]
\centering
\begin{subfigure}[b]{0.75\textwidth}
    \includegraphics[width=\textwidth]{data/fase4/results/sem_sd15_direcciones}
    \caption*{SD~1.5}
\end{subfigure}\\[0.4em]
\begin{subfigure}[b]{0.75\textwidth}
    \includegraphics[width=\textwidth]{data/fase4/results/sem_sd21_direcciones}
    \caption*{SD~2.1}
\end{subfigure}\\[0.4em]
\begin{subfigure}[b]{0.75\textwidth}
    \includegraphics[width=\textwidth]{data/fase4/results/sem_sdxl_direcciones}
    \caption*{SDXL}
\end{subfigure}
\caption{Norma media de las direcciones semánticas por tipo de atributo y modelo.
Las barras de error muestran la desviación típica entre los cinco grupos de cada tipo.}
\label{fig:normas}
\end{figure}

Se observan dos patrones consistentes (Tabla~\ref{tab:normas_dir}, Figura~\ref{fig:normas}).
La jerarquía interna color $>$ iluminación $>$ estilo se mantiene en los tres modelos.
Las normas de SDXL son 2,1--2,8$\times$ las de SD~1.5, un crecimiento proporcionalmente
mayor que el de las normas de los latentes base (factor $\sim 1{,}8\times$), indicando mayor
expresividad semántica neta del doble condicionamiento textual.

\subsection{Consistencia intra-tipo de las direcciones semánticas}
\label{subsec:res:consistencia}

\begin{table}[htb]
\centering
\caption{Similitud coseno media intra-tipo ($\pm$ desviación típica entre pares de grupos).
Valores en $[-1, 1]$; mayor indica mayor universalidad (escena-independencia) de la dirección.}
\label{tab:consistencia}
\small
\begin{tabular}{lrrr}
\toprule
\textbf{Tipo} & \textbf{SD 1.5} & \textbf{SD 2.1} & \textbf{SDXL} \\
\midrule
Color       & $0{,}300 \pm 0{,}061$ & $0{,}229 \pm 0{,}112$ & $\mathbf{0{,}707 \pm 0{,}048}$ \\
Iluminación & $0{,}109 \pm 0{,}091$ & $0{,}155 \pm 0{,}062$ & $\mathbf{0{,}517 \pm 0{,}111}$ \\
Estilo      & $0{,}028 \pm 0{,}081$ & $0{,}010 \pm 0{,}082$ & $\mathbf{0{,}209 \pm 0{,}081}$ \\
\bottomrule
\end{tabular}
\end{table}

\begin{figure}[htbp]
\centering
\begin{subfigure}[b]{0.85\textwidth}
    \includegraphics[width=\textwidth]{data/fase4/results/sem_sd15_angulos}
    \caption*{SD~1.5}
\end{subfigure}\\[0.4em]
\begin{subfigure}[b]{0.85\textwidth}
    \includegraphics[width=\textwidth]{data/fase4/results/sem_sd21_angulos}
    \caption*{SD~2.1}
\end{subfigure}\\[0.4em]
\begin{subfigure}[b]{0.85\textwidth}
    \includegraphics[width=\textwidth]{data/fase4/results/sem_sdxl_angulos}
    \caption*{SDXL}
\end{subfigure}
\caption{Mapa de calor de similitudes coseno entre las quince direcciones semánticas
(cinco por tipo) para cada modelo. En SDXL los cinco grupos de color presentan alta
similitud mutua (bloque rojo superior izquierdo), evidenciando una dirección de color
escena-independiente.}
\label{fig:consistencia}
\end{figure}

SD~1.5 y SD~2.1 muestran consistencias bajas en todos los tipos: las direcciones semánticas
no son universales sino dependientes del contexto de la escena
(Tabla~\ref{tab:consistencia}, Figura~\ref{fig:consistencia}).
SDXL muestra consistencias notablemente más altas: 0,707 para color (ángulo medio de ${\sim}45^{\circ}$
entre pares), 0,517 para iluminación y 0,209 para estilo. En SDXL existe una dirección
universal del color escena-independiente con intensidad equiparable a la del espacio $\mathcal{W}$
de StyleGAN para atributos faciales~\citep{shen2020interfacegan}.

Destaca la ausencia de mejora entre SD~1.5 y SD~2.1 pese al mayor tamaño del codificador
OpenCLIP ViT-H/14: la consistencia de estilo retrocede incluso de 0,028 a 0,010. El salto
cualitativo se produce únicamente con el doble codificador de SDXL.

\subsection{Ortogonalidad inter-tipo (disentanglement)}
\label{subsec:res:ortogonalidad}

\begin{table}[htb]
\centering
\caption{Similitud coseno absoluta media inter-tipo ($\pm$ desviación típica).
Valores en $[0, 1]$; menor indica mejor disentanglement entre tipos de atributo.}
\label{tab:ortogonalidad}
\small
\begin{tabular}{lrrr}
\toprule
\textbf{Par de tipos} & \textbf{SD 1.5} & \textbf{SD 2.1} & \textbf{SDXL} \\
\midrule
Color vs Iluminación  & $0{,}093 \pm 0{,}048$ & $0{,}083 \pm 0{,}052$ & $0{,}119 \pm 0{,}086$ \\
Color vs Estilo       & $0{,}056 \pm 0{,}050$ & $0{,}051 \pm 0{,}039$ & $0{,}227 \pm 0{,}121$ \\
Iluminación vs Estilo & $0{,}068 \pm 0{,}058$ & $0{,}087 \pm 0{,}079$ & $0{,}328 \pm 0{,}130$ \\
\bottomrule
\end{tabular}
\end{table}

Se registra una tensión entre consistencia y disentanglement: SD~1.5 y SD~2.1
presentan excelente ortogonalidad ($|\cos| < 0{,}10$) pero sus direcciones son
contexto-dependientes; SDXL invierte el equilibrio con mayor consistencia pero mayor
entrelazamiento, especialmente entre iluminación y estilo ($|\cos| = 0{,}328$)
(Tabla~\ref{tab:ortogonalidad}).

El entrelazamiento iluminación-estilo en SDXL puede reflejar correlaciones visuales genuinas
del dominio artístico: los estilos plásticos poseen tratamientos lumínicos característicos
(claroscuro barroco, luz difusa impresionista, contraste saturado expresionista), cuya
captura semántica coherente se traduciría necesariamente en correlación entre ambas
direcciones.

\subsection{Alineación de las direcciones semánticas con la base PCA}
\label{subsec:res:energia_pca}

\begin{table}[htb]
\centering
\caption{Fracción media de energía de las direcciones semánticas en los diez primeros
componentes principales ($f_g$, media sobre los cinco grupos del tipo).}
\label{tab:energia_pca}
\small
\begin{tabular}{lrrr}
\toprule
\textbf{Tipo} & \textbf{SD 1.5} & \textbf{SD 2.1} & \textbf{SDXL} \\
\midrule
Color       & 0,72 & 0,67 & \textbf{0,95} \\
Iluminación & 0,56 & 0,50 & \textbf{0,93} \\
Estilo      & 0,51 & 0,51 & \textbf{0,79} \\
\bottomrule
\end{tabular}
\end{table}

En SDXL, el 95\% de la energía de las direcciones de color y el 93\% de las de iluminación
se concentra en los diez primeros componentes principales (Tabla~\ref{tab:energia_pca}).
El PC0 está sistemáticamente alineado con las direcciones de color y el PC3 con las de
iluminación, permitiendo interpretar esas componentes como \textit{dimensiones de color} y
\textit{dimensión de iluminación} del espacio latente de SDXL. Esta propiedad establece
un puente directo con GANSpace~\citep{harkonen2020ganspace}: el PCA sobre $\z_0$ de SDXL
produce componentes semánticamente interpretables sin supervisión, con la ventaja adicional
de que $\z_0$ es directamente accesible mediante inversión DDIM~\citep{song2021ddim} sin
ninguna red de mapeo explícita.

En SD~1.5 y SD~2.1, la fracción de energía (50--72\%) es insuficiente para una alineación
estable; las componentes PCA reflejan principalmente varianza estructural de las escenas,
no organización de atributos visuales. El umbral a partir del cual el PCA descubre semántica
parece situarse entre SD~2.1 y SDXL, coincidiendo con el salto en el condicionamiento textual.

% -----------------------------------------------------------
\section{Fine-tuning y deriva del espacio latente}
\label{sec:res:finetune}
% -----------------------------------------------------------

Los tres modelos se sometieron a fine-tuning LoRA sobre el subconjunto de 500 imágenes de
WikiArt~\citep{saleh2015wikiart}, con los mismos hiperparámetros y precisión, adaptando
únicamente el número de pasos de SDXL (800 frente a 1.000) a las restricciones de
memoria de la GPU T4.
Tras el fine-tuning se regeneró el corpus completo con cada modelo adaptado,
comparando los latentes resultantes ($\z_0^{\text{ft}}$) con los del pretrain ($\z_0^{\text{pre}}$)
para los mismos prompts y semillas. La congelación del VAE garantiza que ambos conjuntos sean
comparables en el mismo espacio vectorial.

\subsection{Deriva geométrica}
\label{subsec:res:deriva_geometrica}

\begin{table}[htb]
\centering
\caption{Métricas de deriva geométrica pre/post fine-tuning WikiArt. $\bar{\delta} =
\|\z_0^\text{ft} - \z_0^\text{pre}\|_2$ medio por latente; $\cos_c$: similitud coseno entre
centroides de la nube pre y post; KS: estadístico del test Kolmogorov--Smirnov sobre las
distribuciones de distancias inter-prompt ($p$ entre paréntesis).}
\label{tab:deriva_geometrica}
\small
\begin{tabular}{lrrrrrrrr}
\toprule
\textbf{Modelo} & \textbf{sep\textsubscript{pre}} & \textbf{sep\textsubscript{ft}} &
$\Delta\text{sep}$ & \textbf{PR\textsubscript{pre}} & \textbf{PR\textsubscript{ft}} &
$\Delta\text{PR}$ & $\cos_c$ & $\bar{\delta}$ \\
\midrule
SD 1.5 & 0,976 & 0,974 & $-0{,}001$ & 12,35 & 11,68 & $-0{,}67$ & 0,991 & \phantom{00}36,6 \\
SD 2.1 & 0,966 & 0,959 & $-0{,}007$ & 12,33 & 10,75 & $-1{,}58$ & 0,993 & \phantom{00}39,6 \\
SDXL   & 1,118 & 1,013 & $\mathbf{-0{,}105}$ & 11,56 & 10,11 & $-1{,}45$ & 0,914 &
$\mathbf{138,4}$ \\
\bottomrule
\end{tabular}
\end{table}

\begin{figure}[htbp]
\centering
\includegraphics[width=\textwidth]{data/fase5/deriva_sd15}
\caption{Deriva latente pre/post fine-tuning WikiArt para SD~1.5. La distribución de
distancias, la separabilidad y el PR apenas varían; $\bar{\delta} = 36{,}6$. La consistencia
semántica y la energía PCA de las direcciones se mantienen estables. La base PCA conserva
una alineación de Grassmann de 0,82 en los diez primeros componentes.}
\label{fig:deriva_sd15}
\end{figure}

\begin{figure}[htbp]
\centering
\includegraphics[width=\textwidth]{data/fase5/deriva_sd21}
\caption{Deriva latente pre/post fine-tuning WikiArt para SD~2.1. La geometría es
prácticamente estable. La consistencia de estilo muestra el único incremento apreciable
entre todos los modelos ($0{,}010 \to 0{,}023$, $+130\%$). La alineación PCA top-10
es la más alta de los tres modelos (0,888).}
\label{fig:deriva_sd21}
\end{figure}

\begin{figure}[htbp]
\centering
\includegraphics[width=\textwidth]{data/fase5/deriva_sdxl}
\caption{Deriva latente pre/post fine-tuning WikiArt para SDXL. El rango del histograma
de desplazamiento por latente alcanza $\sim 250$ frente a $\sim 100$ en SD~1.5 y SD~2.1.
La separabilidad cae de 1,118 a 1,013. La consistencia semántica disminuye en los tres
atributos (color: $-22\%$, iluminación: $-52\%$, estilo: $-37\%$). La alineación PCA
top-10 cae a 0,537.}
\label{fig:deriva_sdxl}
\end{figure}

SD~1.5 y SD~2.1 muestran deriva geométrica prácticamente nula (Tabla~\ref{tab:deriva_geometrica},
Figuras~\ref{fig:deriva_sd15}--\ref{fig:deriva_sd21}).
El ratio de separabilidad varía en menos de 0,01 puntos; los desplazamientos medios por latente
($\bar{\delta} = 36{,}6$ y $39{,}6$) representan el 24\% y el 20\% de las distancias intra-prompt
pretrain, respectivamente. Los centroides se mantienen casi idénticos ($\cos_c > 0{,}99$) y el
test KS no rechaza la igualdad de distribuciones de distancias inter-prompt ($p = 0{,}182$ y
$p = 0{,}665$).

SDXL presenta un perfil cualitativamente distinto (Figura~\ref{fig:deriva_sdxl}).
El desplazamiento medio por latente ($\bar{\delta} = 138{,}4$) equivale al 53\% de la distancia
intra-prompt pretrain ($\mu_\text{intra} \approx 261$): los latentes regenerados de SDXL se han
desplazado más de la mitad de la distancia típica entre semillas del mismo prompt.
El ratio de separabilidad cae 0,105 puntos, perdiendo casi toda la ventaja característica de
SDXL pretrain.
La similitud coseno entre centroides cae a 0,914, registrando la mayor desviación direccional.
A pesar de esta magnitud, el test KS tampoco rechaza la igualdad de distribuciones
($p = 0{,}120$), indicando que la \emph{forma} global de la distribución de distancias se
preserva aunque su posición cambie.

Respecto al participation ratio, SD~2.1 registra la mayor pérdida de dimensionalidad efectiva
($\Delta\text{PR} = -1{,}58$), ligeramente superior a la de SDXL ($-1{,}45$), mientras que
SD~1.5 pierde solo $-0{,}67$ dimensiones. El fine-tuning concentra la varianza en menos
dimensiones en los tres modelos, pero con intensidad mayor en SD~2.1 y SDXL.

\subsection{Rotación de la base PCA tras el fine-tuning}
\label{subsec:res:pca_alineacion}

\begin{table}[htb]
\centering
\caption{Alineación de Grassmann entre los subespacios PCA pretrain y post fine-tuning.
El valor 1 indica subespacios idénticos; valores menores indican rotación de la base.}
\label{tab:pca_alineacion}
\small
\begin{tabular}{lrr}
\toprule
\textbf{Modelo} & \textbf{Alineación top-10} & \textbf{Alineación top-50} \\
\midrule
SD 1.5 & 0,817 & 0,725 \\
SD 2.1 & \textbf{0,888} & \textbf{0,800} \\
SDXL   & 0,537 & 0,282 \\
\bottomrule
\end{tabular}
\end{table}

SD~1.5 y SD~2.1 mantienen alineaciones de 0,82--0,89 en los diez primeros componentes
(Tabla~\ref{tab:pca_alineacion}): el fine-tuning no reorganiza significativamente las
direcciones de mayor varianza.
En SDXL, la alineación cae a 0,537 en los diez primeros componentes y a 0,282 en los
cincuenta primeros: el subespacio de máxima varianza post fine-tuning comparte algo más de
la mitad de la orientación del subespacio pretrain en los componentes principales y menos
de un tercio en el subespacio extendido.

\subsection{Deriva semántica de las direcciones de atributo}
\label{subsec:res:deriva_semantica}

\begin{table}[htb]
\centering
\caption{Consistencia intra-tipo pre y post fine-tuning WikiArt. Los valores entre
paréntesis indican el cambio porcentual relativo.}
\label{tab:deriva_semantica_consistencia}
\small
\begin{tabular}{lcccccc}
\toprule
\textbf{Tipo} & \multicolumn{2}{c}{\textbf{SD 1.5}} & \multicolumn{2}{c}{\textbf{SD 2.1}} &
\multicolumn{2}{c}{\textbf{SDXL}} \\
\cmidrule(lr){2-3}\cmidrule(lr){4-5}\cmidrule(lr){6-7}
 & Pre & Post & Pre & Post & Pre & Post \\
\midrule
Color       & 0,300 & 0,276 {\small($-8\%$)}   & 0,229 & 0,231 {\small($+1\%$)}   &
0,707 & 0,554 {\small($-22\%$)} \\
Iluminación & 0,109 & 0,116 {\small($+6\%$)}   & 0,155 & 0,163 {\small($+5\%$)}   &
0,517 & 0,247 {\small($-52\%$)} \\
Estilo      & 0,028 & 0,027 {\small($-4\%$)}   & 0,010 & 0,023 {\small($+130\%$)} &
0,209 & 0,132 {\small($-37\%$)} \\
\bottomrule
\end{tabular}
\end{table}

SD~1.5 y SD~2.1 exhiben estabilidad semántica casi total: los cambios en la consistencia
intra-tipo son inferiores al 8\% relativo en todos los atributos, con la única excepción de la
consistencia de estilo en SD~2.1, que pasa de 0,010 a 0,023 ($+130\%$ relativo)
(Tabla~\ref{tab:deriva_semantica_consistencia}).
Este incremento, aunque grande en términos porcentuales, parte de un valor de base prácticamente
nulo; el resultado post fine-tuning (0,023) sigue siendo inferior en un orden de magnitud a la
consistencia de estilo pretrain de SDXL (0,209).

SDXL muestra una deriva semántica severa en todos los atributos: $-22\%$ en color, $-52\%$
en iluminación y $-37\%$ en estilo. Los valores post fine-tuning de SDXL siguen siendo
superiores a los pretrain de SD~1.5 y SD~2.1, confirmando que la representación semántica
de SDXL, aunque perturbada, mantiene ventaja cualitativa. La disminución más severa
corresponde a la iluminación, el atributo cuya dirección estaba parcialmente alineada con
el estilo antes del fine-tuning ($|\cos| = 0{,}328$, Tabla~\ref{tab:ortogonalidad}), sugiriendo
que la reorganización inducida por WikiArt afecta especialmente a las propiedades con
correlación visual en el dominio artístico.

\subsection{Deriva en la magnitud y la energía PCA de las direcciones semánticas}
\label{subsec:res:deriva_normas_energia}

Además de la consistencia, el fine-tuning modifica la magnitud de las direcciones y su
alineación con la base de varianza principal. La Tabla~\ref{tab:deriva_normas_energia}
resume la norma media y la energía PCA pre y post fine-tuning para cada tipo de atributo.

\begin{table}[htb]
\centering
\caption{Norma media y energía PCA ($f_g$ en los diez primeros componentes) de las
direcciones semánticas antes y después del fine-tuning WikiArt. Los cambios porcentuales
relativos se indican entre paréntesis. Solo se muestra SDXL para los atributos con cambio
apreciable; en SD~1.5 y SD~2.1 todos los deltas son inferiores al 7\% y se omiten.}
\label{tab:deriva_normas_energia}
\small
\begin{tabular}{lccccccc}
\toprule
 & \multicolumn{3}{c}{\textbf{Norma media ($\norm{\mathbf{d}_g}$)}} &
\multicolumn{3}{c}{\textbf{Energía PCA ($f_g$)}} \\
\cmidrule(lr){2-4}\cmidrule(lr){5-7}
\textbf{Tipo} & \textbf{Pre} & \textbf{Post} & \textbf{$\Delta$\%} &
\textbf{Pre} & \textbf{Post} & \textbf{$\Delta$\%} \\
\midrule
\multicolumn{7}{l}{\textit{SDXL}} \\
Color       & 183,4 & 129,4 & $-29\%$ & 0,948 & 0,925 & $-2\%$ \\
Iluminación & 149,8 & \phantom{1}98,7 & $-34\%$ & 0,932 & 0,845 & $-9\%$ \\
Estilo      & 106,5 & \phantom{1}90,9 & $-15\%$ & 0,794 & 0,789 & $-1\%$ \\
\addlinespace
\multicolumn{7}{l}{\textit{SD~2.1} (referencia)} \\
Estilo      & \phantom{1}64,8 & \phantom{1}60,2 & $-7\%$ & 0,511 & 0,553 & $+8\%$ \\
\bottomrule
\end{tabular}
\end{table}

En SDXL, las normas de las direcciones de color e iluminación caen un 29\% y un 34\%
respectivamente (Tabla~\ref{tab:deriva_normas_energia}). Esta reducción de magnitud acompaña
y amplifica la caída de consistencia: no solo las direcciones se vuelven menos universales,
sino que la diferencia vectorial entre los dos valores de cada atributo se acorta. La energía PCA
se mantiene relativamente estable (caída de 2--9\%), indicando que las direcciones post
fine-tuning de SDXL, aunque más cortas y menos consistentes, siguen proyectándose
principalmente sobre el subespacio de máxima varianza.

En SD~2.1, la norma de la dirección de estilo cae un 7\% pero su energía PCA mejora
ligeramente ($+8\%$), coherente con el pequeño incremento de consistencia (+130\% relativo)
observado en la Tabla~\ref{tab:deriva_semantica_consistencia}: el fine-tuning sobre WikiArt
da cierto orden semántico a la dirección de estilo de SD~2.1, aunque en magnitud
absolutamente insuficiente para alcanzar los valores de SDXL.

\subsection{Comparativa cross-modelo de la deriva latente}
\label{subsec:res:deriva_cross}

\begin{figure}[htbp]
\centering
\includegraphics[width=\textwidth]{data/deriva_cross_model}
\caption{Comparativa cross-modelo de la deriva latente tras fine-tuning WikiArt.
Fila superior (izq.\ a der.): $\Delta$sep, $\Delta$PR y alineación PCA top-10.
Fila inferior: $\Delta$consistencia intra-tipo, $\Delta$energía PCA de las direcciones
semánticas y desplazamiento latente medio $\bar{\delta}$.
El modelo SDXL (verde) domina todas las métricas de deriva geométrica y semántica.}
\label{fig:deriva_cross}
\end{figure}

La Figura~\ref{fig:deriva_cross} sintetiza la comparativa entre los tres modelos.
El patrón es uniforme y consistente: SD~1.5 y SD~2.1 presentan deltas cercanos a cero en
prácticamente todas las métricas, mientras que SDXL exhibe las mayores magnitudes de cambio:
\begin{itemize}[noitemsep, leftmargin=2em]
    \item Separabilidad: $\Delta\text{sep} = -0{,}105$ (frente a $-0{,}001$ y $-0{,}007$)
    \item Desplazamiento latente: $\bar{\delta} = 138{,}4$ ($\times3{,}5$--$3{,}8$ respecto a SD~1.5/SD~2.1)
    \item Alineación PCA top-10: 0,537 (frente a 0,817 y 0,888)
    \item $\Delta$consistencia color: $-0{,}153$; $\Delta$consistencia iluminación: $-0{,}270$
\end{itemize}

La notable excepción es $\Delta\text{PR}$: SD~2.1 registra la mayor pérdida de dimensionalidad
efectiva ($-1{,}58$), ligeramente superior a la de SDXL ($-1{,}45$), mientras que SD~1.5
pierde solo $-0{,}67$. Este patrón sugiere que SD~2.1, aunque geométricamente estable a nivel
de distancias, experimenta la mayor concentración dimensional de varianza tras el fine-tuning:
el espacio latente mantiene su posición pero se vuelve localmente más plano.

% -----------------------------------------------------------
\section*{Síntesis del capítulo}
% -----------------------------------------------------------

\noindent
Los resultados se articulan en torno a tres observaciones centrales.

\medskip
\noindent PR$\approx 12$: convergencia geométrica universal.
La dimensionalidad intrínseca es prácticamente idéntica en los tres modelos con independencia
de sus dimensiones nominales, confirmando que esta propiedad está fijada por el proceso de
difusión, no por la arquitectura.

\medskip
\noindent Consistencia semántica: divergencia dependiente del codificador.
El doble codificador de SDXL induce consistencias de atributo 2,3--21 veces superiores a las
de SD~1.5 y SD~2.1. El umbral de universalidad semántica en $\z_0$ no se cruza escalando un
único codificador CLIP.

\medskip
\noindent Fine-tuning WikiArt: inversión del orden de fragilidad.
SDXL, con la representación pretrain más organizada, sufre la mayor perturbación geométrica
y semántica; SD~1.5 y SD~2.1 muestran deriva prácticamente nula. La organización semántica
pretrain amplifica la susceptibilidad a la disrupción por adaptación de dominio.

% ===========================================================
%  REFERENCIAS (solo las citadas en este capítulo)
% ===========================================================


% ===========================================================
%  APÉNDICES
% ===========================================================
\appendix

\label{app:prompts}

\section{Prompts base}

\begin{longtable}{clp{9.5cm}}
\caption{Lista completa de los 20 prompts base del corpus (diseñados con Claude \citep{anthropic2024claude}).}
\label{tab:prompts_completos} \\
\toprule
\textbf{ID} & \textbf{Dominio} & \textbf{Prompt} \\
\midrule
\endfirsthead
\toprule \textbf{ID} & \textbf{Dominio} & \textbf{Prompt} \\ \midrule
\endhead
\midrule \multicolumn{3}{r}{\small\textit{continúa\ldots}} \endfoot
\bottomrule \endlastfoot
base\_01 & Naturaleza muerta & \textit{a red apple on a wooden table, studio lighting, photorealistic} \\
base\_02 & Animal / exterior & \textit{a golden retriever dog sitting in a sunlit park} \\
base\_03 & Urbano / noche & \textit{a futuristic city skyline at night with neon lights} \\
base\_04 & Paisaje natural & \textit{a mountain lake at sunrise with mist and reflections} \\
base\_05 & Interior & \textit{a minimalist white living room with afternoon light} \\
base\_06 & Arquitectura & \textit{an ancient stone castle on a hilltop, overcast sky} \\
base\_07 & Tropical / costa & \textit{a tropical beach with turquoise water and palm trees} \\
base\_08 & Bosque / otoño & \textit{a forest path in autumn with golden leaves} \\
base\_09 & Bodegón bebida & \textit{a cup of coffee with latte art, warm morning light} \\
base\_10 & Tecnología / cielo & \textit{a rocket launching at dusk, long exposure photography} \\
base\_11 & Gastronomía & \textit{a bowl of ramen with steam rising, food photography} \\
base\_12 & Animal / interior & \textit{a black cat on a windowsill with rain outside} \\
base\_13 & Objetos / urbano & \textit{a vintage red bicycle against a textured brick wall} \\
base\_14 & Naturaleza / agua & \textit{a waterfall in a lush tropical jungle} \\
base\_15 & Fauna / sabana & \textit{an elephant in the African savanna at golden hour} \\
base\_16 & Arquitectura religiosa & \textit{Gothic cathedral interior with stained glass windows} \\
base\_17 & Montaña / noche & \textit{a snow-covered mountain peak under a starry night sky} \\
base\_18 & Cabaña / invierno & \textit{a wooden cabin in a snowy pine forest} \\
base\_19 & Campo / verano & \textit{a sunflower field under a vivid blue sky} \\
base\_20 & Arquitectura moderna & \textit{a modern glass skyscraper reflecting blue sky} \\
\end{longtable}

\section{Prompts de atributo}

\begin{longtable}{llp{4.4cm}p{4.4cm}}
\caption{Lista completa de los 30 prompts de atributo (15 grupos $\times$ 2 valores, diseñados con Claude \citep{anthropic2024claude}).}
\label{tab:prompts_atributo} \\
\toprule
\textbf{Tipo} & \textbf{Grupo} & \textbf{Valor 1} & \textbf{Valor 2} \\
\midrule
\endfirsthead
\toprule \textbf{Tipo} & \textbf{Grupo} & \textbf{Valor 1} & \textbf{Valor 2} \\ \midrule
\endhead
\midrule \multicolumn{4}{r}{\small\textit{continúa\ldots}} \endfoot
\bottomrule \endlastfoot
Color & color\_apple & \textit{a warm red apple on a white surface} & \textit{a cool blue apple on a white surface} \\
Color & color\_car & \textit{a warm orange sports car on the road} & \textit{a cool blue sports car on the road} \\
Color & color\_umbrella & \textit{a warm yellow umbrella in urban rain} & \textit{a cool teal umbrella in urban rain} \\
Color & color\_flower & \textit{a warm red rose in a vase} & \textit{a cool blue rose in a vase} \\
Color & color\_balloon & \textit{a warm orange hot air balloon} & \textit{a cool blue hot air balloon} \\
Ilum. & light\_portrait & \textit{portrait, bright high-key lighting} & \textit{portrait, dark low-key lighting} \\
Ilum. & light\_kitchen & \textit{modern kitchen, bright morning light} & \textit{modern kitchen, dark night LED light} \\
Ilum. & light\_forest & \textit{forest path, bright midday sunlight} & \textit{forest path, dark sunset light} \\
Ilum. & light\_street & \textit{cobblestone street, bright sunny afternoon} & \textit{cobblestone street, dark night with lamppost} \\
Ilum. & light\_still\_life & \textit{fruit still life, bright studio lighting} & \textit{fruit still life, single candle light} \\
Estilo & style\_mountain & \textit{snowy mountain, photorealistic} & \textit{snowy mountain, impressionist painting} \\
Estilo & style\_portrait & \textit{young man portrait, photorealistic} & \textit{young man portrait, watercolor style} \\
Estilo & style\_city & \textit{aerial city view, photorealistic} & \textit{aerial city view, digital art style} \\
Estilo & style\_forest & \textit{light through trees, photorealistic} & \textit{light through trees, fantasy painting} \\
Estilo & style\_cat & \textit{cat on sofa, photorealistic} & \textit{cat on sofa, Studio Ghibli style} \\
\end{longtable}

% ===========================================================
%  REFERENCIAS BIBLIOGRÁFICAS
% ===========================================================
\renewcommand{\bibname}{Referencias bibliográficas}
\addcontentsline{toc}{chapter}{Referencias bibliográficas}
\begin{thebibliography}{99}

\bibitem[Anthropic(2024)]{anthropic2024claude}
Anthropic. (2024).
\textit{Claude} [modelo de lenguaje grande].
\url{https://www.anthropic.com}

\bibitem[Betker et al.(2023)]{betker2023dalle3}
Betker, J., et al.\ (2023).
\textit{Improving image generation with better captions}.
OpenAI.

\bibitem[Black Forest Labs(2024)]{bfl2024flux}
Black Forest Labs. (2024).
\textit{Flux.1}.
\url{https://github.com/black-forest-labs/flux}

\bibitem[Dinh et al.(2017)]{dinh2017realnvp}
Dinh, L., Sohl-Dickstein, J., y Bengio, S. (2017).
Density estimation using real-valued non-volume preserving transformations.
\textit{International Conference on Learning Representations}.

\bibitem[Efros y Leung(1999)]{efros1999texture}
Efros, A.~A., y Leung, T.~K. (1999).
Texture synthesis by non-parametric sampling.
\textit{Proceedings of the 7th IEEE International Conference on Computer Vision (ICCV)}, 1033--1038.

\bibitem[Esser et al.(2024)]{esser2024sd3}
Esser, P., et al.\ (2024).
Scaling rectified flow transformers for high-resolution image synthesis.
\textit{Proceedings of the 41st International Conference on Machine Learning (ICML)}.

\bibitem[Facco et al.(2017)]{facco2017twonn}
Facco, E., d'Errico, M., Rodriguez, A., y Laio, A. (2017).
Estimating the intrinsic dimension of datasets by a minimal neighborhood information.
\textit{Scientific Reports, 7}(1), 12140.

\bibitem[Gal et al.(2022)]{gal2022textual}
Gal, R., et al.\ (2022).
An image is worth one word: Personalizing text-to-image generation using textual inversion.
\textit{International Conference on Learning Representations 2023}.

\bibitem[Gatys et al.(2016)]{gatys2016style}
Gatys, L.~A., Ecker, A.~S., y Bethge, M. (2016).
Image style transfer using convolutional neural networks.
\textit{IEEE/CVF Conference on Computer Vision and Pattern Recognition (CVPR)}, 2414--2423.

\bibitem[Google LLC(2017)]{google2017colab}
Google LLC. (2017).
\textit{Google Colaboratory}.
\url{https://colab.research.google.com}

\bibitem[Goodfellow et al.(2014)]{goodfellow2014gans}
Goodfellow, I., et al.\ (2014).
Generative adversarial nets.
\textit{Advances in Neural Information Processing Systems (NeurIPS), 27}.

\bibitem[H\"{a}rkonen et al.(2020)]{harkonen2020ganspace}
H\"{a}rkonen, E., Hertzmann, A., Lehtinen, J., y Paris, S. (2020).
GANSpace: Discovering interpretable GAN controls.
\textit{Advances in Neural Information Processing Systems (NeurIPS), 33}.

\bibitem[Hertz et al.(2022)]{hertz2022p2p}
Hertz, A., et al.\ (2022).
Prompt-to-Prompt image editing with cross attention control.
\textit{International Conference on Learning Representations 2023}.

\bibitem[Ho et al.(2020)]{ho2020ddpm}
Ho, J., Jain, A., y Abbeel, P. (2020).
Denoising diffusion probabilistic models.
\textit{Advances in Neural Information Processing Systems (NeurIPS), 33}, 6840--6851.

\bibitem[Ho y Salimans(2022)]{ho2022cfg}
Ho, J., y Salimans, T. (2022).
Classifier-free diffusion guidance.
\textit{NeurIPS Workshop on Deep Generative Models and Downstream Applications}.

\bibitem[Hu et al.(2022)]{hu2022lora}
Hu, E.~J., Shen, Y., Wallis, P., Allen-Zhu, Z., Li, Y., Wang, S., Wang, L., y Chen, W. (2022).
LoRA: Low-rank adaptation of large language models.
\textit{International Conference on Learning Representations}.

\bibitem[Ilharco et al.(2021)]{ilharco2021openclip}
Ilharco, G., et al.\ (2021).
\textit{OpenCLIP}.
Zenodo. \url{https://doi.org/10.5281/zenodo.5143773}

\bibitem[Karras et al.(2019)]{karras2019stylegan}
Karras, T., Laine, S., y Aila, T. (2019).
A style-based generator architecture for generative adversarial networks.
\textit{IEEE/CVF Conference on Computer Vision and Pattern Recognition (CVPR)}, 4401--4410.

\bibitem[Karras et al.(2020)]{karras2020stylegan2}
Karras, T., Laine, S., Aittala, M., Hellsten, J., Lehtinen, J., y Aila, T. (2020).
Analyzing and improving the image quality of StyleGAN.
\textit{IEEE/CVF Conference on Computer Vision and Pattern Recognition (CVPR)}, 8110--8119.

\bibitem[Kingma y Dhariwal(2018)]{kingma2018glow}
Kingma, D.~P., y Dhariwal, P. (2018).
Glow: Generative flow with invertible $1{\times}1$ convolutions.
\textit{Advances in Neural Information Processing Systems (NeurIPS), 31}.

\bibitem[Kingma y Welling(2013)]{kingma2013vae}
Kingma, D.~P., y Welling, M. (2013).
Auto-encoding variational Bayes.
\textit{International Conference on Learning Representations 2014}.

\bibitem[Kirkpatrick et al.(2017)]{kirkpatrick2017ewc}
Kirkpatrick, J., et al.\ (2017).
Overcoming catastrophic forgetting in neural networks.
\textit{Proceedings of the National Academy of Sciences, 114}(13), 3521--3526.

\bibitem[Krizhevsky et al.(2012)]{krizhevsky2012alexnet}
Krizhevsky, A., Sutskever, I., y Hinton, G.~E. (2012).
ImageNet classification with deep convolutional neural networks.
\textit{Advances in Neural Information Processing Systems (NeurIPS), 25}.

\bibitem[Kwon et al.(2023)]{kwon2023diffusion}
Kwon, M., Jeong, J., y Uh, Y. (2023).
Diffusion models already have a semantic latent space.
\textit{International Conference on Learning Representations}.

\bibitem[Lipman et al.(2022)]{lipman2022flowmatching}
Lipman, Y., et al.\ (2022).
Flow matching for generative modeling.
\textit{International Conference on Learning Representations 2023}.

\bibitem[Liu et al.(2022)]{liu2022rectified}
Liu, X., Gong, C., y Liu, Q. (2022).
Flow straight and fast: Learning to generate and transfer data with rectified flow.
\textit{International Conference on Learning Representations 2023}.

\bibitem[McCloskey y Cohen(1989)]{mccloskey1989catastrophic}
McCloskey, M., y Cohen, N.~J. (1989).
Catastrophic interference in connectionist networks: The sequential learning problem.
\textit{Psychology of Learning and Motivation, 24}, 109--165.

\bibitem[Mokady et al.(2023)]{mokady2023nulltext}
Mokady, R., et al.\ (2023).
Null-text inversion for editing real images using guided diffusion models.
\textit{IEEE/CVF Conference on Computer Vision and Pattern Recognition (CVPR)}.

\bibitem[Peebles y Xie(2023)]{peebles2023dit}
Peebles, W., y Xie, S. (2023).
Scalable diffusion models with transformers.
\textit{IEEE/CVF International Conference on Computer Vision (ICCV)}.

\bibitem[Perlin(1985)]{perlin1985noise}
Perlin, K. (1985).
An image synthesizer.
\textit{ACM SIGGRAPH Computer Graphics, 19}(3), 287--296.

\bibitem[Podell et al.(2023)]{podell2023sdxl}
Podell, D., English, Z., Lacey, K., Blattmann, A., Dockhorn, T., M\"{u}ller, J.,
Penna, J., y Rombach, R. (2023).
SDXL: Improving latent diffusion models for high-resolution image synthesis.
\textit{arXiv}. \url{https://arxiv.org/abs/2307.01952}

\bibitem[Pope et al.(2021)]{pope2021intrinsic}
Pope, P., Zhu, C., Abdelkader, A., Goldblum, M., y Goldstein, T. (2021).
The intrinsic dimension of images and its impact on learning.
\textit{International Conference on Learning Representations}.

\bibitem[Radford et al.(2015)]{radford2015dcgan}
Radford, A., Metz, L., y Chintala, S. (2015).
Unsupervised representation learning with deep convolutional generative adversarial networks.
\textit{arXiv}:1511.06434.

\bibitem[Radford et al.(2021)]{radford2021clip}
Radford, A., et al.\ (2021).
Learning transferable visual models from natural language supervision.
\textit{Proceedings of the 38th International Conference on Machine Learning (ICML)}, 8748--8763.

\bibitem[Rombach et al.(2022)]{rombach2022ldm}
Rombach, R., Blattmann, A., Lorenz, D., Esser, P., y Ommer, B. (2022).
High-resolution image synthesis with latent diffusion models.
\textit{IEEE/CVF Conference on Computer Vision and Pattern Recognition (CVPR)}, 10684--10695.

\bibitem[Ruiz et al.(2023)]{ruiz2023dreambooth}
Ruiz, N., et al.\ (2023).
DreamBooth: Fine tuning text-to-image diffusion models for subject-driven generation.
\textit{IEEE/CVF Conference on Computer Vision and Pattern Recognition (CVPR)}.

\bibitem[Saleh y Elgammal(2015)]{saleh2015wikiart}
Saleh, B., y Elgammal, A. (2015).
Large-scale classification of fine-art paintings: Learning the right metric on the right feature.
\textit{arXiv}. \url{https://arxiv.org/abs/1505.00855}

\bibitem[Schuhmann et al.(2022)]{schuhmann2022laion}
Schuhmann, C., et al.\ (2022).
LAION-5B: An open large-scale dataset for training next generation image-text models.
\textit{Advances in Neural Information Processing Systems (NeurIPS), 35}.

\bibitem[Shen et al.(2020)]{shen2020interfacegan}
Shen, Y., Gu, J., Tang, X., y Zhou, B. (2020).
Interpreting the latent space of GANs for semantic face editing.
\textit{IEEE/CVF Conference on Computer Vision and Pattern Recognition (CVPR)}, 9243--9252.

\bibitem[Shen y Zhou(2021)]{shen2021sefa}
Shen, Y., y Zhou, B. (2021).
Closed-form factorization of latent semantics in GANs.
\textit{IEEE/CVF Conference on Computer Vision and Pattern Recognition (CVPR)}.

\bibitem[Song et al.(2021)]{song2021ddim}
Song, J., Meng, C., y Ermon, S. (2021).
Denoising diffusion implicit models.
\textit{International Conference on Learning Representations}.

\bibitem[Stability AI(2022)]{stabilityai2022sd2}
Stability AI. (2022).
\textit{Stable Diffusion 2}.
\url{https://github.com/Stability-AI/stablediffusion}

\bibitem[Tang et al.(2023)]{tang2023emergent}
Tang, L., et al.\ (2023).
Emergent correspondence from image diffusion.
\textit{Advances in Neural Information Processing Systems (NeurIPS), 36}.

\bibitem[van den Oord et al.(2016)]{vandeoord2016pixelrnn}
van den Oord, A., Kalchbrenner, N., y Kavukcuoglu, K. (2016).
Pixel recurrent neural networks.
\textit{Proceedings of the 33rd International Conference on Machine Learning (ICML)}.

\bibitem[Wang et al.(2004)]{wang2004ssim}
Wang, Z., Bovik, A.~C., Sheikh, H.~R., y Simoncelli, E.~P. (2004).
Image quality assessment: From error visibility to structural similarity.
\textit{IEEE Transactions on Image Processing, 13}(4), 600--612.

\bibitem[Wu et al.(2021)]{wu2021stylespace}
Wu, Z., Lischinski, D., y Shechtman, E. (2021).
StyleSpace analysis: Disentangled controls for StyleGAN image generation.
\textit{IEEE/CVF Conference on Computer Vision and Pattern Recognition (CVPR)}.

\end{thebibliography}

\end{document}
